\section{Discussion}\label{sec:discussion}

\subsection{What the Results Support}

In the setting of Section~\ref{sec:host}, v2 is the fastest of the route matchers measured here in
\HOneSuperior{} of \HOneCells{} cells, and not slower in one more. ``Fastest'' is said of a cell only
against the competitors that take part in it. The competitors left out are named with their reasons
in Section~\ref{sec:run}, and the paper does not claim that they are slower.

The advantage grows with the table. v2's work per lookup barely depends on $m$. On \shape{static} it
ran \HTwoVtwoNetStaticTen{} net instructions per lookup at $m=\SizeTen$ and
\HTwoVtwoNetStaticHundredThousand{} at $m=\SizeHundredThousand$, while the fastest competitor's count
rose from \HTwoFastestNetStaticTen{} to \HTwoFastestNetStaticHundredThousand. A walk over path
segments takes one step per segment, so its length follows the depth of the route, not the number of
routes. The width of a node matters only through its scan of up to \VTwoLinearEdges{} literal
children, or one probe of their hashed slots. In every cell, v2 also took fewer branch misses than
the median competitor.

\subsection{Small Tables Against gin}\label{sec:gin}

Every loss or partial result of v2 in the lookup comparison sits at $m=\SizeTen$ and against gin.
These are the \HTwoLost{} H2 losses, the one cell that does not pass H1, and the one cell that passes
non-inferiority only. With ten routes, first words rarely repeat, so most nodes below the root have
one child. A node of gin's tree compares its whole prefix at once, while v2 takes one step per
segment. The development round attributed the instruction loss at \shape{mixed-disjoint} to this
difference. The rework tried to remove it with chains, and its rules rejected every attempt
(Section~\ref{sec:rework}).

An instruction count is not a time, however. At \shape{param-last} with $m=\SizeTen$, v2 ran more net
instructions than gin and was still faster, with the ratio \HOneRatioParamLastTen. Part of v2's count
is the harness's adapter, which copies the captured values into the harness's own structure and folds
them. gin folds its values inside Go, and its null pass does not. Appendix~\ref{app:lookup} lists
these differences between the arms. We kept them, because removing them would change the harness for
every arm.

At \shape{wild} with $m=\SizeTen$, the development builds favoured v2. The run differs from those
builds in the instruction set that every arm targets and in the build level of the Go arms. Both are
candidate causes of the change in this cell, and neither was tested.

\subsection{Build Time and the Compile-Time Table}

The rework cut v2's build to \AbBuildMin{} to \AbBuildMax{} of the pre-release's time and reduced its
heap. It was not enough in \HFourLostBuild{} cells, at \shape{static} and \shape{param-first} with
$m=\SizeTenThousand$ and \SizeHundredThousand. We did not profile what remains of the build there. A
server builds its table once at start, and v2's slowest build of the largest tables took
\HFourBuildMaxMs~ms. As pre-specified, H4's claim is not made.

The compile-time table passed its per-cell test in \HSevenHeld{} of \HSevenCells{} cells. The rule
needs all of them, so H7 does not hold. The one cell that fails has \HSevenLostAboveMargin{} pairs at
or above the margin, where the rule allows \HSevenAllowedAbove, and its interval lies wholly above
one. One of them, at table seed \HSevenLostSeed, ran more cycles with the same instructions. The
rework had already given both tables one layout. What remains between them includes where the table
lies, in the binary or on the heap, and where the code lies. The pre-specified run with a heap copy,
made afterwards, found single slow pairs at other seeds and in each arm. Neither the table seed nor
the table's place in memory explains the cell, and its cause stays open. Zero-overhead abstraction is therefore not established
here for compile-time tables~\cite{stroustrup2012foundations}.

The costs set a practical limit. A compile-time table of \SizeThousand{} routes took
\ClangThousandCompileMin{} to \ClangThousandCompileMax~s to compile with clang. With MSVC it took
\MsvcPeakMinMiB{} to \MsvcPeakMaxMiB~MiB of memory, and MSVC crashed under the smallest budget of its
search. A compile-time table is cheap at \SizeTen{} and \SizeHundred{} routes and expensive at
\SizeThousand.

\subsection{End to End}

The lookup difference against v1 reached the client in every cell of H6. In the minimal server the
null router's time per request was about \HSixNullUsMin~$\mu$s, so routing is a visible share of a
request in a server this lean. The measured ratio stayed within \HSixPredMin{} to \HSixPredMax{} of
the ratio that the lookup times predict. The share would be smaller in a full framework with more
work per request.

\subsection{Threats to Validity}\label{sec:threats}

\begin{itemize}
\item \textbf{One processor model and one core.} Every time comes from one core of one \HostLArch{}
processor. Another microarchitecture may rank close cells differently, above all the cells at
$m=\SizeTen$ against gin.
\item \textbf{One compiler per language.} C and C++ were compiled by clang only, Rust by rustc and Go by
its own toolchain. GCC and MSVC builds were not timed. MSVC appears only in the costs of compile-time
tables.
\item \textbf{The native instruction set.} Every arm was built for L's own processor. A generic build
would give other code and other times, which may matter most for the small tables against gin.
\item \textbf{Loopback for H6.} The end-to-end run used loopback, one server core and one request in
flight per connection. A network, more cores or pipelining would change the share of routing.
\item \textbf{The minimal server.} It parses, routes and answers a fixed body. A full framework adds
work per request, which dilutes the difference between routers.
\item \textbf{Development on the same host.} The engineering of v2 ran on L, on other seeds. The run's
seeds were held out, but the design choices may still suit L's microarchitecture better than others.
\item \textbf{The Drogon transcription.} Drogon's router needs a running application, so its lookup was
transcribed into the harness from its source. The file's hash and line ranges are pinned. A
transcription can still differ from the router as Drogon runs it.
\item \textbf{The gin shim.} gin's tree is vendored unchanged, but the parts of gin's engine that a
lookup uses are reproduced in a small shim. They are the reset of the context, the method's tree and
the search in gin's default configuration. chi's arm uses a similar shim of its multiplexer. A shim
can differ from the engine it stands for.
\item \textbf{The repeated ring.} Every timed pass repeats the same \RingSize{} queries. Predictors and
caches learn the ring, and a large table's hot part is at most the ring's routes. Real traffic, with
more distinct paths, may favour other structures. The Zipf and miss-heavy rings explore two departures
only.
\item \textbf{Foreign-function boundaries.} The Rust arms are called through one C function per lookup,
and the Go arms run their timed loops inside Go, entered through cgo. Their nulls remove the loops, not
every cost of the boundary.
\item \textbf{The cut for slow arms.} A competitor whose passes ran on a prefix of the ring has
optimistic times and counters. None was the fastest competitor of an H1 cell, but the cut competitors
enter the median of H2's branch misses, which Table~\ref{tab:htwo} gives with and without them.
\end{itemize}
